Let $u=x+y$. The second equation gives $(y-x)u=2$, so $u\ne0$. The first equation and $n>0$ force $u>0$. Also $y-x=\dfrac{2}{u}$, and therefore
\[
n=\dfrac{x^2+y^2}{u}=\dfrac{u}{2}+\dfrac{2}{u^3}.
\]
Each $u>0$ determines exactly one pair $x=\dfrac{1}{2}(u-\dfrac{2}{u})$, $y=\dfrac{1}{2}(u+\dfrac{2}{u})$.
The function on the right tends to infinity at both ends of $(0,\infty)$, decreases until $u^4=12$, and increases afterward, since its derivative is $\dfrac{1}{2}-\dfrac{6}{u^4}$. Exactly one solution occurs at its minimum. There $n=\dfrac{2u}{3}$, giving
\[
n^4=\dfrac{16}{81}\cdot12=\dfrac{256}{27}.
\]
Thus $\boxed{p+q=283}$.
